\section{总结与延伸}

本讲最重要的启发是：科研 Agent 的竞争力来自\textbf{系统编排}，而非单一模型分数。Virtual Lab 说明多角色协作可以提升开放性科研任务的探索效率；Paper2Agent 说明知识传播可以从静态文档升级为可调用能力；Agents for Science 说明我们可以用会议级数据衡量人机协作与 AI 评审的真实表现。

\subsection*{全讲总结表}

\begin{table}[H]
\centering
\small
\begin{tabular}{p{0.20\textwidth}p{0.33\textwidth}p{0.35\textwidth}}
\toprule
主题 & 关键结论 & 对实践者的直接动作 \\
\midrule
Virtual Lab 组织范式 & PI 路由 + 多专家协作 + Critic 审查可显著提升开放任务稳定性 & 先定义角色与会议协议，再优化 prompt；把 Critic 设为默认角色 \\
Agent School 与可靠性 & memory/sandbox/HITL 是长期可用性的基础，不是可选项 & 建立外部记忆库、可重放执行环境、关键节点人工闸门 \\
Paper2Agent 与 MCP & 先构建 robust MCP 再执行任务，准确性和效率更优 & 对高价值论文先做环境重建与接口封装，沉淀为复用资产 \\
Agent-to-Agent 协作 & 协议驱动协作可发现跨论文新组合机会 & 建立方法 Agent 与数据 Agent 的匹配与对接流程 \\
Agents for Science 证据 & AI 评审可发现细节错误，但存在模型偏置与谄媚风险 & 使用多模型评审委员会 + 人类抽检 + 校准策略 \\
治理与边界 & 科研 Agent 需要可审计、可核验、可问责的责任结构 & 将引用核验、复现检查、human sign-off 纳入默认流程 \\
\bottomrule
\end{tabular}
\end{table}

\subsection*{进一步阅读}

\begin{itemize}
    \item Swanson et al., Virtual Lab（课程中介绍为 Nature 2025 相关工作，建议结合公开视频与论文正文阅读）。
    \item Miao et al., Paper2Agent（课程中介绍为 arXiv 2025 相关工作，关注 MCP 封装与评测设置）。
    \item Agents for Science 会议资料与公开评审（关注 AI involvement checklist 与引用核验机制）。
    \item Model Context Protocol (MCP) 官方文档（理解资源暴露、工具调用和协议边界）。
    \item 关于 AI 评审偏置与 sycophancy 的近期研究（用于设计多模型评审和评分校准）。
\end{itemize}

从课程落地角度看，最值得先做的小型试点是：选择一个复现成本高但影响大的论文子领域，先落地 Paper2Agent；再把经过验证的能力接入团队版 Virtual Lab，最后用 checklist 记录人机分工与错误类型，形成自己的``小型 Agents for Science''数据闭环。

\end{document}
